
\begin{center}\Large\bfseries Дополнительные материалы к статье о СМО с отрицательными заявками\end{center}
\section*{1. Назначение и исходные данные}
Приложение сопровождает версию 8 статьи «Расчёт многоканальных СМО с отрицательными заявками». Оно содержит дополнительные характеристики, сетки и параметры проверки. Основным таблицам соответствуют 96 имитационных прогонов. Дополнительные серии имеют отдельно заданные параметры и объём наблюдения. Везде ДИ означает полуширину приближённого 95\%-го интервала, а \(\Delta=100(\text{расчёт}-\text{ИМ})/\text{ИМ}\). Все числа сформированы из сохранённых JSON-результатов, без ручного округления промежуточных моментов.

\noindent TABLECAPX Таблица SEQXtab_1X. Стандартные отклонения основных 12 примеров. Имитация совпадает с основной статьёй.\par

\setlength{\tabcolsep}{2pt}
\begin{tabular}{p{4.16cm}p{1.61cm}p{2.83cm}p{1.28cm}p{1.61cm}p{2.83cm}p{1.28cm}}
\toprule
 & \multicolumn{3}{c}{$\sigma_W$} & \multicolumn{3}{c}{$\sigma_V$} \\
Вариант & Расчёт & ИМ $\pm$ ДИ & $\Delta$, \% & Расчёт & ИМ $\pm$ ДИ & $\Delta$, \% \\
\midrule
\multicolumn{7}{l}{Коэффициент вариации: 0,8} \\
RCS, удаление & 0,9129 & 0,9340 ± 0,0230 & -2,26 & 1,6972 & 1,7093 ± 0,0157 & -0,70 \\
Очищение & 0,5519 & 0,5555 ± 0,0079 & -0,65 & 1,2635 & 1,2652 ± 0,0052 & -0,13 \\
RCS, новая длительность & 1,9618 & 1,9601 ± 0,0489 & 0,09 & 2,6671 & 2,6657 ± 0,0347 & 0,05 \\
Все каналы, новая длительность & 2,7947 & 2,7967 ± 0,0779 & -0,07 & 3,4332 & 3,4350 ± 0,0638 & -0,05 \\
RCS, сохранённая длительность & 1,7606 & 1,7460 ± 0,0447 & 0,84 & 2,4855 & 2,4742 ± 0,0342 & 0,46 \\
Все каналы, сохранённая длительность & 2,1566 & 2,1149 ± 0,0474 & 1,97 & 2,8568 & 2,8203 ± 0,0399 & 1,29 \\
\multicolumn{7}{l}{Коэффициент вариации: 1,2} \\
RCS, удаление & 1,0268 & 1,0252 ± 0,0268 & 0,16 & 2,0866 & 2,0891 ± 0,0159 & -0,12 \\
Очищение & 0,5699 & 0,5696 ± 0,0086 & 0,05 & 1,4337 & 1,4333 ± 0,0040 & 0,03 \\
RCS, новая длительность & 2,0208 & 1,9836 ± 0,0444 & 1,87 & 3,0849 & 3,0572 ± 0,0319 & 0,91 \\
Все каналы, новая длительность & 1,5746 & 1,5817 ± 0,0274 & -0,45 & 2,6559 & 2,6569 ± 0,0181 & -0,04 \\
RCS, сохранённая длительность & 2,7736 & 2,8138 ± 0,0911 & -1,43 & 3,8523 & 3,8838 ± 0,0700 & -0,81 \\
Все каналы, сохранённая длительность & 3,8830 & 3,8867 ± 0,2009 & -0,09 & 4,8665 & 4,8699 ± 0,1703 & -0,07 \\
\bottomrule
\end{tabular}




\section*{2. Уточнение сохранённой работы}
В основных примерах \(J=4\), \(N=100\); для RCS \(H=24\). Экстраполяция \(K=8,12,16\) по \(1/K\) имеет веса \(2,-9,8\) и применяется к начальным моментам. Таблица конечных сеток позволяет отделить сами расчёты от экстраполированного результата. Малые отличия при увеличении \(N,H\) включают остаточную ошибку линейного решателя. Увеличение числа состояний само по себе не является доказательством сходимости стационарных моментов.

\noindent TABLECAPX Таблица SEQXtab_2X. Конечные сетки: вариант и \(CV\), число фаз на класс, начальные моменты ожидания, время расчёта в секундах.\par

\setlength{\tabcolsep}{2pt}
\begin{tabular}{p{4.44cm}p{1.13cm}p{3.31cm}p{3.31cm}p{2.46cm}}
\toprule
Вариант, $CV$ & $K$ & $w_1$ & $w_2$ & Время, с \\
\midrule
RCS, 0,8 & 4 & 1,239397 & 6,322977 & 7,6 \\
RCS, 0,8 & 8 & 1,139711 & 5,206458 & 120,1 \\
RCS, 0,8 & 12 & 1,105919 & 4,853345 & 83,0 \\
RCS, 0,8 & 16 & 1,088892 & 4,680347 & 337,2 \\
RCS, 1,2 & 4 & 1,864266 & 15,055481 & 8,1 \\
RCS, 1,2 & 8 & 1,726325 & 12,552254 & 124,3 \\
RCS, 1,2 & 12 & 1,679069 & 11,750098 & 126,2 \\
RCS, 1,2 & 16 & 1,655167 & 11,355380 & 520,7 \\
Все, 0,8 & 4 & 1,421020 & 8,207528 & 1,4 \\
Все, 0,8 & 8 & 1,347865 & 7,216742 & 40,5 \\
Все, 0,8 & 12 & 1,322678 & 6,894780 & 178,7 \\
Все, 0,8 & 16 & 1,309931 & 6,735760 & 1146,4 \\
Все, 1,2 & 4 & 2,324342 & 23,043622 & 10,3 \\
Все, 1,2 & 8 & 2,267271 & 21,460974 & 173,0 \\
Все, 1,2 & 12 & 2,247481 & 20,944201 & 115,2 \\
Все, 1,2 & 16 & 2,237492 & 20,690638 & 1357,7 \\
\bottomrule
\end{tabular}




\noindent TABLECAPX Таблица SEQXtab_3X. Экстраполированные средние при сохранённой работе и дополнительная серия имитационного моделирования.\par

\setlength{\tabcolsep}{2pt}
\begin{tabular}{p{4.16cm}p{1.61cm}p{2.83cm}p{1.28cm}p{1.61cm}p{2.83cm}p{1.28cm}}
\toprule
 & \multicolumn{3}{c}{$w_1$} & \multicolumn{3}{c}{$v_1$} \\
Вариант & Расчёт & ИМ $\pm$ ДИ & $\Delta$, \% & Расчёт & ИМ $\pm$ ДИ & $\Delta$, \% \\
\midrule
RCS, CV=0,8 & 1,0373 & 1,0457 ± 0,0082 & -0,81 & 3,1680 & 3,1782 ± 0,0100 & -0,32 \\
RCS, CV=1,2 & 1,5824 & 1,5888 ± 0,0203 & -0,40 & 3,7280 & 3,7345 ± 0,0214 & -0,18 \\
Все каналы, CV=0,8 & 1,2711 & 1,2654 ± 0,0168 & 0,45 & 3,4458 & 3,4409 ± 0,0190 & 0,14 \\
Все каналы, CV=1,2 & 2,2071 & 2,2292 ± 0,0503 & -0,99 & 4,4209 & 4,4443 ± 0,0533 & -0,53 \\
\bottomrule
\end{tabular}




При увеличении стационарного предела общих прерываний с \(N=100\) до 160 на сетке \(CV=1{,}2,J=4,K=16\) среднее ожидание меняется на \(6{,}00\cdot10^{-6}\), второй начальный момент --- на \(6{,}57\cdot10^{-4}\) (около 0{,}0032\%). Для RCS увеличение \(N\) или \(H\) на контрольной сетке меняет первые два момента менее чем на \(1{,}3\cdot10^{-6}\). Это изменения конечных сеток; веса экстраполяции могут усиливать остаточную ошибку. Результаты обеих точностей линейного решателя сохранены отдельно.

Сопоставление более грубой тройки фаз, числа классов и другого аргумента экстраполяции:

\noindent TABLECAPX Таблица SEQXtab_4X. Относительные изменения, \%, относительно основной экстраполяции. Они не являются строгими границами ошибки.\par

\setlength{\tabcolsep}{2pt}
\begin{tabular}{p{4.44cm}p{2.79cm}p{2.79cm}p{2.79cm}p{2.79cm}}
\toprule
Вариант, $CV$ & $w_1$: грубее & $w_1$: $J=3$ & $w_1$: $1/\sqrt K$ & $\sigma_W$: $1/\sqrt K$ \\
\midrule
RCS, 0,8 & 0,0191 & -0,1228 & -0,2201 & 1,0739 \\
RCS, 1,2 & 0,0177 & -0,4003 & -0,3029 & 0,8178 \\
Все, 0,8 & 0,0019 & -0,0854 & -0,2109 & 0,4408 \\
Все, 1,2 & -0,0176 & -0,2377 & -0,0752 & 0,3036 \\
\bottomrule
\end{tabular}




\noindent FIGBLOCKX0X\par

\noindent PANELCAPX0X0X (a) Конечные сетки и экстраполяции.\par

\noindent MAINCAPX0X Рисунок SEQXfig_1X. Сходимость среднего ожидания при сохранённой работе. Полосы показывают доверительные интервалы дополнительной серии имитационного моделирования. Коэффициенты экстраполяции не используют имитационные значения.\par
\section*{3. Расширенный эксперимент}
В каждом условии \(b_1=2{,}1\), \(\lambda=\rho n/b_1\). Исходный закон при \(CV=0{,}8\) --- гамма-распределение, при \(CV=1{,}2\) --- положительный \(H_2\), согласованный с первыми тремя моментами гамма-распределения. Для каждого условия выполнены 8 прогонов с исключением первых 50\,000 уходов из статистики и наблюдением следующих 500\,000 уходов. При сохранённой работе основная сетка \(J=3\), \(K=4,8,12\); для двух каналов \(N=120,H=32\), для трёх \(N=240,H=40\). Для общих прерываний свёртка с временем собственной работы устраняет необходимость в \(H\). Средняя сетка более нагруженных случаев дополнительно проверена при \(N=400\).

\noindent TABLECAPX Таблица SEQXtab_5X. Дополнительные условия. Номер сохраняется во всех следующих таблицах.\par

\setlength{\tabcolsep}{2pt}
\begin{tabular}{p{0.76cm}p{7.65cm}p{1.04cm}p{1.89cm}p{1.89cm}p{1.89cm}}
\toprule
№ & Механизм & $n$ & $\rho$ & $CV$ & $\delta$ \\
\midrule
1 & RCS, удаление & 2 & 0,45 & 0,8 & 0,30 \\
2 & RCS, удаление & 3 & 0,85 & 1,2 & 0,60 \\
3 & Очищение & 2 & 0,90 & 0,8 & 0,15 \\
4 & Очищение & 4 & 1,20 & 1,2 & 0,30 \\
5 & RCS, новая длительность & 2 & 0,50 & 0,8 & 0,60 \\
6 & RCS, новая длительность & 4 & 0,85 & 1,2 & 0,60 \\
7 & Все каналы, новая длительность & 2 & 0,40 & 0,8 & 0,60 \\
8 & Все каналы, новая длительность & 3 & 0,85 & 1,2 & 0,60 \\
9 & RCS, сохранённая длительность & 2 & 0,50 & 0,8 & 0,04 \\
10 & RCS, сохранённая длительность & 3 & 0,85 & 1,2 & 0,04 \\
11 & Все каналы, сохранённая длительность & 2 & 0,50 & 0,8 & 0,04 \\
12 & Все каналы, сохранённая длительность & 3 & 0,80 & 1,2 & 0,04 \\
\bottomrule
\end{tabular}




\noindent TABLECAPX Таблица SEQXtab_6X. Средние в дополнительных условиях. Для гамма-распределения в вариантах удаления и новой длительности использовано 9 фаз.\par

\setlength{\tabcolsep}{2pt}
\begin{tabular}{p{4.16cm}p{1.61cm}p{2.83cm}p{1.28cm}p{1.61cm}p{2.83cm}p{1.28cm}}
\toprule
 & \multicolumn{3}{c}{$w_1$} & \multicolumn{3}{c}{$v_1$} \\
Вариант & Расчёт & ИМ $\pm$ ДИ & $\Delta$, \% & Расчёт & ИМ $\pm$ ДИ & $\Delta$, \% \\
\midrule
1 & 0,1909 & 0,1904 ± 0,0011 & 0,28 & 1,7005 & 1,6999 ± 0,0027 & 0,04 \\
2 & 0,3195 & 0,3191 ± 0,0048 & 0,13 & 1,5240 & 1,5243 ± 0,0071 & -0,02 \\
3 & 0,9133 & 0,9139 ± 0,0071 & -0,06 & 2,3485 & 2,3493 ± 0,0072 & -0,03 \\
4 & 0,4585 & 0,4607 ± 0,0074 & -0,46 & 1,4992 & 1,5025 ± 0,0086 & -0,22 \\
5 & 0,9254 & 0,9214 ± 0,0061 & 0,44 & 3,3072 & 3,3018 ± 0,0078 & 0,16 \\
6 & 1,6984 & 1,7218 ± 0,0254 & -1,36 & 3,6576 & 3,6827 ± 0,0262 & -0,68 \\
7 & 0,7021 & 0,7015 ± 0,0058 & 0,08 & 3,2342 & 3,2340 ± 0,0081 & 0,01 \\
8 & 1,2073 & 1,1999 ± 0,0047 & 0,62 & 2,9820 & 2,9751 ± 0,0060 & 0,23 \\
9 & 0,7065 & 0,7143 ± 0,0052 & -1,09 & 2,9037 & 2,9140 ± 0,0063 & -0,35 \\
10 & 6,2171 & 6,2476 ± 0,1060 & -0,49 & 8,3978 & 8,4282 ± 0,1076 & -0,36 \\
11 & 0,8813 & 0,8795 ± 0,0093 & 0,20 & 3,1358 & 3,1334 ± 0,0109 & 0,08 \\
12 & 9,9564 & 9,9801 ± 0,1662 & -0,24 & 12,2979 & 12,3195 ± 0,1668 & -0,18 \\
\bottomrule
\end{tabular}




\noindent TABLECAPX Таблица SEQXtab_7X. Стандартные отклонения в дополнительных условиях. Выходы за интервалы не исключались из результатов.\par

\setlength{\tabcolsep}{2pt}
\begin{tabular}{p{4.16cm}p{1.61cm}p{2.83cm}p{1.28cm}p{1.61cm}p{2.83cm}p{1.28cm}}
\toprule
 & \multicolumn{3}{c}{$\sigma_W$} & \multicolumn{3}{c}{$\sigma_V$} \\
Вариант & Расчёт & ИМ $\pm$ ДИ & $\Delta$, \% & Расчёт & ИМ $\pm$ ДИ & $\Delta$, \% \\
\midrule
1 & 0,6100 & 0,6083 ± 0,0034 & 0,28 & 1,4385 & 1,4379 ± 0,0025 & 0,04 \\
2 & 0,8413 & 0,8399 ± 0,0150 & 0,17 & 1,6618 & 1,6621 ± 0,0123 & -0,02 \\
3 & 1,5760 & 1,5787 ± 0,0158 & -0,17 & 2,0138 & 2,0151 ± 0,0129 & -0,07 \\
4 & 1,0022 & 1,0060 ± 0,0157 & -0,38 & 1,5797 & 1,5834 ± 0,0115 & -0,23 \\
5 & 1,8110 & 1,8031 ± 0,0118 & 0,44 & 2,7302 & 2,7232 ± 0,0100 & 0,26 \\
6 & 2,6597 & 2,6896 ± 0,0602 & -1,11 & 3,5004 & 3,5260 ± 0,0485 & -0,72 \\
7 & 1,6145 & 1,6189 ± 0,0159 & -0,27 & 2,7742 & 2,7770 ± 0,0111 & -0,10 \\
8 & 2,0527 & 2,0386 ± 0,0138 & 0,69 & 2,8229 & 2,8138 ± 0,0129 & 0,32 \\
9 & 1,5134 & 1,5221 ± 0,0105 & -0,57 & 2,4459 & 2,4511 ± 0,0091 & -0,21 \\
10 & 7,9318 & 8,1138 ± 0,3531 & -2,24 & 8,4091 & 8,5808 ± 0,3372 & -2,00 \\
11 & 1,9304 & 1,9217 ± 0,0255 & 0,45 & 2,8434 & 2,8382 ± 0,0208 & 0,18 \\
12 & 13,1308 & 12,9447 ± 0,4779 & 1,44 & 13,5749 & 13,3929 ± 0,4657 & 1,36 \\
\bottomrule
\end{tabular}




\noindent TABLECAPX Таблица SEQXtab_8X. Стоимость исходных сеток расширенного эксперимента и максимальное абсолютное изменение начальных моментов при проверке отсечений. Для гамма-распределения время положительного PH-уточнения сохранено отдельно в исходных данных.\par

\setlength{\tabcolsep}{2pt}
\begin{tabular}{p{0.94cm}p{4.25cm}p{4.72cm}p{5.67cm}}
\toprule
№ & Время сеток, с & Пиковая память, МиБ & Изменение при уточнении \\
\midrule
1 & 1,8 & 95,1 & 0,000000 \\
2 & 1,0 & 107,8 & 0,000000 \\
3 & 1,7 & 95,9 & 0,000000 \\
4 & 1,2 & 115,4 & 0,000000 \\
5 & 1,3 & 101,0 & 0,000000 \\
6 & 1,5 & 116,2 & 0,000000 \\
7 & 1,3 & 101,8 & 0,000000 \\
8 & 1,0 & 107,7 & 0,000000 \\
9 & 8,6 & 331,5 & — \\
10 & 1579,9 & 367,9 & 0,001762 \\
11 & 6,4 & 99,5 & — \\
12 & 791,0 & 333,4 & 0,006041 \\
\bottomrule
\end{tabular}




Память указана по максимальному объёму оперативной памяти, занятой отдельным процессом (RSS); время --- по полному времени выполнения расчёта, при одновременной работе нескольких процессов. Прочерк означает, что отдельное увеличение предела для данного условия не выполнялось; нулевое значение после округления означает малое измеренное изменение. Выбор прямого или итерационного решателя зависит от размера системы, поэтому время не обязано монотонно расти с пределом очереди. Проверка половин прогонов и закона Литтла сохранена для каждого условия. Отдельный 95\%-й интервал не гарантирует одновременного покрытия всех строк и характеристик.

В условии 8 средние ожидания и пребывания, а также стандартное отклонение ожидания находятся вне интервалов; двухфазный расчёт даёт те же значения. В условии 9 вне интервалов находятся оба средних. Чувствительность результатов условия 9 к числу классов и фаз приведена в следующей таблице. Все приведённые ниже средние остаются ниже соответствующих интервалов: \(0{,}71427\pm0{,}00520\) и \(2{,}91397\pm0{,}00629\).

\noindent TABLECAPX Таблица SEQXtab_9X. Дополнительная чувствительность условия 9 по числу классов и фаз. Экстраполируются начальные моменты по \(1/K\).\par

\setlength{\tabcolsep}{2pt}
\begin{tabular}{p{1.42cm}p{4.72cm}p{4.72cm}p{4.72cm}}
\toprule
$J$ & $K$ & $w_1$ & $v_1$ \\
\midrule
3 & 4, 8, 12 & 0,706511 & 2,903677 \\
3 & 8, 12, 16 & 0,706355 & 2,903541 \\
4 & 4, 8, 12 & 0,707690 & 2,904901 \\
4 & 8, 12, 16 & 0,707573 & 2,904799 \\
5 & 4, 8, 12 & 0,708162 & 2,905364 \\
5 & 8, 12, 16 & 0,708080 & 2,905292 \\
\bottomrule
\end{tabular}






\section*{4. Положительные фазовые распределения}
Для гамма-распределения при \(CV=0{,}8\) использованы квадратуры с \(J=2,4,8\), соответствующие 3, 5 и 9 фазам. На контрольной сетке аргумента максимальное отклонение функции распределения от исходного гамма-распределения составляет соответственно 0{,}00633; 0{,}000808 и 0{,}0000954. Это дискретная проверка, не доказанная равномерная граница. Первые \(2J-1\) моментов сохраняются по точности квадратуры.

\noindent TABLECAPX Таблица SEQXtab_10X. Среднее ожидание при 3, 5 и 9 положительных фазах и изменение между 5 и 9 фазами, \%.\par

\setlength{\tabcolsep}{2pt}
\begin{tabular}{p{5.76cm}p{2.55cm}p{2.55cm}p{2.55cm}p{2.17cm}}
\toprule
Механизм & 3 фазы & 5 фаз & 9 фаз & $\Delta$, \% \\
\midrule
RCS, удаление & 0,409683 & 0,409965 & 0,409967 & 0,000550 \\
Очищение & 0,191780 & 0,192061 & 0,192062 & 0,000527 \\
RCS, новая длительность & 1,205480 & 1,205963 & 1,205966 & 0,000272 \\
Все каналы, новая длительность & 1,860551 & 1,859758 & 1,859769 & 0,000572 \\
\bottomrule
\end{tabular}




\noindent TABLECAPX Таблица SEQXtab_11X. Распределение Эрланга четвёртого порядка: \(CV=0{,}5\), \(n=3\), \(\lambda=0{,}6\), \(\delta=0{,}3\), новая длительность. Отдельные 8 прогонов на вариант.\par

\setlength{\tabcolsep}{2pt}
\begin{tabular}{p{4.16cm}p{1.61cm}p{2.83cm}p{1.28cm}p{1.61cm}p{2.83cm}p{1.28cm}}
\toprule
 & \multicolumn{3}{c}{$w_1$} & \multicolumn{3}{c}{$v_1$} \\
Вариант & Расчёт & ИМ $\pm$ ДИ & $\Delta$, \% & Расчёт & ИМ $\pm$ ДИ & $\Delta$, \% \\
\midrule
RCS & 0,1992 & 0,1993 ± 0,0016 & -0,06 & 2,5824 & 2,5826 ± 0,0013 & -0,01 \\
Все каналы & 0,4335 & 0,4326 ± 0,0037 & 0,22 & 3,0838 & 3,0829 ± 0,0056 & 0,03 \\
\bottomrule
\end{tabular}




Для рационального двухфазного ПЛС при \(CV=0{,}8\) найдено положительное представление Кокса:


\[\boldsymbol\alpha=(1,0),\]

\[S=\begin{pmatrix}-0{,}674119662&0{,}511518043\\0&-1{,}230642243\end{pmatrix}.\]


Вероятность перехода во вторую фазу равна 0{,}758794120. Эквивалентность представлений установлена по ПЛС и проверена на четырёх вариантах очереди. Полученное двухфазное распределение приближает исходное гамма-распределение.

\section*{5. Воспроизводимость и границы проверки}
Исходные JSON-файлы содержат неокруглённые начальные моменты, параметры, отсечения, время и память. Файлы отдельных прогонов содержат все начальные состояния генераторов случайных чисел. Для RCS с сохранённой работой используется база 202611072000 с шагом 100 между значениями \(CV\); для 12 дополнительных условий --- 202611073000 с шагом 100 по условию; для распределения Эрланга четвёртого порядка --- 202611075000 с шагом 100 по варианту. Внутри каждой серии 8 прогонов используются смещения 0--7. Прогоны не исключались по результату.

Уравнения и событийный имитатор расположены в сопровождающем каталоге расчётов. Полный список команд приведён в README. В комплект включены фазовые расчёты и результаты изменения численных параметров. Контроль 49 тестами охватывает известные частные случаи и независимые реализации; он не устанавливает равномерную точность для всех нагрузок и законов обслуживания. Экстраполяции на сохранённой работе требуют отдельной оценки при иных параметрах.

